\documentclass[12pt,a4paper]{article}

\usepackage[a4paper,margin=1in]{geometry}
\usepackage{graphicx}
\usepackage{booktabs}
\usepackage{amsmath,amssymb}
\usepackage{float}
\usepackage{hyperref}
\usepackage{xcolor}
\usepackage{listings}
\usepackage{enumitem}
\usepackage{fancyhdr}
\usepackage{longtable}

\hypersetup{colorlinks=true,linkcolor=blue,urlcolor=blue}

\pagestyle{fancy}
\fancyhf{}
\lhead{ICS1512}
\rhead{Machine Learning Algorithms Laboratory}
\cfoot{\thepage}

\lstset{
language=Python,
basicstyle=\ttfamily\small,
numbers=left,
frame=single,
breaklines=true,
showstringspaces=false,
keywordstyle=\color{blue},
commentstyle=\color{gray},
stringstyle=\color{purple}
}

\begin{document}

\begin{tabular}{|p{4cm}|p{11cm}|}
\hline
Experiment No. & 6 \\ \hline
Experiment Title & Bagging, Boosting, and Stacked Ensemble Models
\footnote{\textbf{GitHub Repository: }\url{https://github.com/Nightingales21/ICS1512-Machine-Learning-Labratory}}\\ \hline
Student Name & Nitin \\ \hline
Register Number & 312224700142 \\ \hline
Faculty & Dr. Poreddy Ajay Kumar Reddy \\ \hline
Submission Date & \today \\ \hline
\end{tabular}

\section{Objective}
\begin{itemize}
\item To understand ensemble learning strategies: Bagging, Boosting, and Stacking.
\item To implement Bagging and Boosting classifiers.
\item To build a Stacked Ensemble model using multiple base learners.
\item To compare ensemble models in terms of accuracy, stability, and generalization.
\item To analyze the effect of ensemble methods on bias and variance.
\end{itemize}

\section{Dataset Description}
\begin{longtable}{|p{6cm}|p{8cm}|}
\hline
Dataset Name & Wisconsin Diagnostic Breast Cancer (WDBC) \\ \hline
Dataset Source & UCI Machine Learning Repository (loaded via
\texttt{sklearn.datasets.load\_breast\_cancer}, identical to Experiment 5) \\ \hline
Number of Samples & 569 \\ \hline
Number of Features & 30 numerical attributes \\ \hline
Number of Classes & 2 (Malignant, Benign) \\ \hline
Missing Values & 0 (verified programmatically) \\ \hline
Train--Test Split & 80\% / 20\%, stratified, \texttt{random\_state=42} \\ \hline
\end{longtable}

Class balance: 357 benign samples (62.7\%) and 212 malignant samples
(37.3\%). As in Experiment 5, scikit-learn's loader encodes
\texttt{target=0} as malignant and \texttt{target=1} as benign; this
encoding is used consistently in all confusion matrices below. This
experiment builds directly on Experiment 5's Decision Tree baseline, now
extending it into three ensemble families.

\section{Preprocessing Steps}
\begin{itemize}
\item \textbf{Class label encoding:} handled identically to Experiment 5
(0 = malignant, 1 = benign).
\item \textbf{Missing value handling:} None required; zero missing cells.
\item \textbf{Feature scaling:} Not required for Bagging, AdaBoost, or
Gradient Boosting, since all use decision-tree base learners which are
scale-invariant. The Stacked Ensemble's SVM base learner is, in principle,
scale-sensitive, but was left unscaled here to keep all three ensemble
families directly comparable on identical input features; this is noted as
a possible refinement in the Discussion.
\item \textbf{Train--test split:} \texttt{train\_test\_split} with
\texttt{test\_size=0.2}, \texttt{stratify=y}, \texttt{random\_state=42}.
\end{itemize}

\section{Exploratory Data Analysis}
The reusable function \texttt{generate\_eda\_summary()} from
\texttt{ml\_lab\_utils.py} (Experiment 1, unmodified) was used to produce
the consolidated 12-panel EDA figure below; it is identical in structure to
Experiment 5's EDA since the same dataset is used.

\begin{figure}[H]
\centering
\includegraphics[width=\linewidth]{figures/eda_breast_cancer.png}
\caption{Consolidated 12-panel EDA summary for the Wisconsin Breast Cancer dataset.}
\label{fig:eda}
\end{figure}
\textbf{Inference:} As established in Experiment 5, the dataset has no
missing values, a moderate 62.7\%/37.3\% class imbalance, and strong
multicollinearity among the size-related features (\texttt{radius},
\texttt{perimeter}, \texttt{area}). This redundancy is precisely the kind of
structure ensemble methods can exploit differently: Bagging's random
feature subsampling and Boosting's sequential re-weighting both implicitly
handle correlated features differently than a single greedy tree, which is
explored quantitatively in Section~10 (bias--variance analysis).

\section{Implementation Details}

\textbf{Bagging (Bootstrap Aggregation)} trains multiple Decision Tree base
estimators independently, each on a different bootstrap-resampled subset of
the training data (and optionally a random feature subset), then averages
their predictions. Because each tree sees different data, their individual
errors are partially independent, so averaging reduces the ensemble's
overall \emph{variance} without materially increasing bias.

\textbf{Boosting} trains base estimators \emph{sequentially}, with each new
estimator focusing on the samples the previous ones misclassified.
\textbf{AdaBoost} re-weights misclassified samples so subsequent weak
learners pay more attention to them; \textbf{Gradient Boosting} fits each
new estimator to the residual errors of the current ensemble. Because each
step explicitly targets the current model's mistakes, boosting primarily
reduces \emph{bias}, allowing an ensemble of simple (high-bias) learners to
collectively fit complex decision boundaries.

\textbf{Stacked Ensemble} trains several heterogeneous base learners (here:
SVM, Gaussian Naive Bayes, and Decision Tree) and a meta-learner (Logistic
Regression) that learns how to optimally combine the base learners'
predictions, using cross-validated out-of-fold predictions to avoid
information leakage into the meta-learner's training data.

\begin{lstlisting}[language=Python, caption={Baseline Bagging via the reusable classification function, from \texttt{ml\_lab\_utils.py} (Experiment 1).}]
baseline_models = {"Bagging (Baseline)": BaggingClassifier(
    estimator=DecisionTreeClassifier(random_state=42), random_state=42)}
baseline_results_df, baseline_fitted = train_evaluate_classification(
    baseline_models, X_train, X_test, y_train, y_test, scale=False)
\end{lstlisting}

\begin{lstlisting}[language=Python, caption={Stacked Ensemble construction, from \texttt{experiment6\_ensemble\_breast\_cancer.py}.}]
base_learners = [
    ("svm", SVC(probability=True, random_state=42)),
    ("nb", GaussianNB()),
    ("dt", DecisionTreeClassifier(random_state=42)),
]
stacking_model = StackingClassifier(
    estimators=base_learners,
    final_estimator=LogisticRegression(max_iter=2000),
    cv=StratifiedKFold(5, shuffle=True, random_state=42))
\end{lstlisting}

\section{Experimental Setup}
\begin{tabular}{ll}
\toprule
Python Version & 3.12 \\
Scikit-Learn Version & 1.8.0 \\
IDE & Jupyter Notebook / VS Code \\
Operating System & Ubuntu 24.04 LTS \\
Hardware & Standard lab workstation (CPU only) \\
\bottomrule
\end{tabular}

\vspace{0.3cm}
\textbf{Environment activation command (as mandated by the lab manual):}
\begin{lstlisting}[language=bash]
source /opt/anaconda3/bin/activate anaconda-navigation
\end{lstlisting}

\section{Hyperparameter Evaluation Results}

\textbf{Table 1: Bagging Hyperparameter Evaluation}
\begin{tabular}{ccc c}
\toprule
n\_estimators & max\_samples & Avg CV Accuracy (\%) & Avg CV F1 Score \\
\midrule
10  & 0.5 & 95.16 & 0.9610 \\
10  & 0.7 & 95.16 & 0.9611 \\
10  & 1.0 & 95.60 & 0.9648 \\
50  & 0.5 & \textbf{96.04} & 0.9681 \\
50  & 0.7 & 95.82 & 0.9665 \\
50  & 1.0 & 96.04 & 0.9683 \\
100 & 0.5 & 95.38 & 0.9630 \\
100 & 0.7 & 96.04 & 0.9683 \\
100 & 1.0 & 96.04 & \textbf{0.9684} \\
\bottomrule
\end{tabular}

\vspace{0.2cm}
\noindent \textbf{Best Bagging (full 3-parameter GridSearchCV):}
n\_estimators=50, max\_samples=0.5, max\_features=0.5, CV accuracy = 96.04\%.

\vspace{0.3cm}
\textbf{Table 2: Boosting (AdaBoost) Hyperparameter Evaluation}
\begin{tabular}{ccc c}
\toprule
n\_estimators & learning\_rate & Avg CV Accuracy (\%) & Avg CV F1 Score \\
\midrule
50  & 0.01 & 92.53 & 0.9407 \\
50  & 0.10 & 95.38 & 0.9633 \\
50  & 1.00 & 96.70 & 0.9740 \\
100 & 0.01 & 93.41 & 0.9480 \\
100 & 0.10 & 96.04 & 0.9686 \\
100 & 1.00 & \textbf{98.02} & \textbf{0.9844} \\
200 & 0.01 & 93.63 & 0.9496 \\
200 & 0.10 & 97.14 & 0.9773 \\
200 & 1.00 & 97.80 & 0.9826 \\
\bottomrule
\end{tabular}

\vspace{0.2cm}
\noindent \textbf{Best AdaBoost:} n\_estimators=100, learning\_rate=1.0, CV
accuracy = 98.02\%. \textbf{Best Gradient Boosting} (separate 3-parameter
grid over n\_estimators, learning\_rate, max\_depth): n\_estimators=200,
learning\_rate=0.1, max\_depth=2, CV accuracy = 96.70\%. AdaBoost was
therefore selected as the representative ``Boosting'' model for all
subsequent comparisons.

\vspace{0.3cm}
\textbf{Table 3: Stacked Ensemble Evaluation}
\begin{tabular}{llcc}
\toprule
Base Models & Meta Learner & Avg CV Accuracy (\%) & Avg CV F1 Score \\
\midrule
SVM+NB+DT & Logistic Regression & \textbf{95.38} & \textbf{0.9635} \\
SVM+NB+DT & Decision Tree (depth 3) & 94.07 & 0.9528 \\
NB+DT     & Logistic Regression & 94.29 & 0.9551 \\
\bottomrule
\end{tabular}

\vspace{0.2cm}
\noindent \textbf{Best Stacked Ensemble configuration:} base models
SVM+Naive Bayes+Decision Tree with a Logistic Regression meta-learner,
CV accuracy = 95.38\%. A Logistic Regression meta-learner outperformed a
Decision Tree meta-learner, and dropping SVM from the base-learner set
(NB+DT only) reduced accuracy, confirming SVM's margin-based decision
boundary contributes complementary information the other two base learners
lack.

\vspace{0.3cm}
\textbf{Hyperparameter Tuning Results Summary}
\begin{tabular}{lllc}
\toprule
Model & Search Method & Best Parameters & Best CV Accuracy \\
\midrule
Bagging            & GridSearchCV (5-fold) & n\_est=50, samples=0.5, feat=0.5 & 0.9604 \\
AdaBoost           & GridSearchCV (5-fold) & n\_est=100, lr=1.0              & \textbf{0.9802} \\
Gradient Boosting  & GridSearchCV (5-fold) & n\_est=200, lr=0.1, depth=2     & 0.9670 \\
Stacked Ensemble   & Manual combo search   & SVM+NB+DT / LogReg              & 0.9538 \\
\bottomrule
\end{tabular}

\section{Performance Comparison of Ensemble Models}

\textbf{Table 4: Performance Comparison of Ensemble Models (test set)}
\begin{tabular}{lcccc}
\toprule
Model & Accuracy (\%) & Precision & Recall & F1 Score \\
\midrule
Bagging           & 95.61 & 0.9561 & 0.9561 & 0.9560 \\
Boosting (AdaBoost) & 95.61 & 0.9569 & 0.9561 & 0.9558 \\
Stacked Ensemble  & 95.61 & 0.9569 & 0.9561 & 0.9558 \\
\bottomrule
\end{tabular}

\vspace{0.2cm}
\noindent All three tuned ensembles reach an identical 95.61\% accuracy on
this particular 20\% test split (a tie at this sample size, 114 test
points, corresponds to differing by at most 1--2 predictions). Their
5-fold cross-validated accuracies, shown next, tell a more differentiated
and reliable story.

\begin{figure}[H]
\centering
\includegraphics[width=\linewidth]{figures/ensemble_comparison_bar.png}
\caption{Ensemble model comparison: Accuracy, Precision, Recall, F1 (test set).}
\label{fig:bar}
\end{figure}

\section{5-Fold Cross-Validation Comparison}

\begin{tabular}{lcccccc}
\toprule
Model & Fold 1 & Fold 2 & Fold 3 & Fold 4 & Fold 5 & Average \\
\midrule
Bagging          & 0.9560 & 0.9560 & 0.9341 & 0.9670 & 0.9890 & 0.9604 \\
Boosting         & 0.9890 & 1.0000 & 0.9451 & 0.9780 & 0.9890 & \textbf{0.9802} \\
Stacked Ensemble & 0.9341 & 0.9780 & 0.9341 & 0.9341 & 0.9890 & 0.9538 \\
\bottomrule
\end{tabular}

\begin{figure}[H]
\centering
\includegraphics[width=0.8\linewidth]{figures/cv_accuracy_comparison.png}
\caption{5-fold cross-validation accuracy: Bagging vs. Boosting vs. Stacked Ensemble.}
\label{fig:cv}
\end{figure}
\textbf{Inference:} Although all three models tied on the single test
split, boosting's 5-fold CV average (98.02\%) is clearly the highest and
even reaches a perfect 1.000 on Fold 2, while the stacked ensemble is the
weakest performer on cross-validation despite combining three different
model families -- illustrating that more heterogeneous components do not
automatically guarantee the best accuracy if the meta-learner cannot fully
exploit the base learners' complementary strengths (in this case, a simple
linear meta-learner over three already-strong-but-correlated classifiers
does not clearly beat a well-tuned sequential boosting ensemble).

\vspace{0.3cm}
\textbf{Stability Comparison (CV standard deviation, lower = more stable)}
\begin{tabular}{lcc}
\toprule
Model & CV Mean Accuracy & CV Std Dev \\
\midrule
Bagging          & 0.9604 & \textbf{0.0179} \\
Boosting         & 0.9802 & 0.0189 \\
Stacked Ensemble & 0.9538 & 0.0245 \\
\bottomrule
\end{tabular}
\quad Bagging is the most stable (lowest fold-to-fold variance) of the
three, consistent with its theoretical purpose of variance reduction; the
Stacked Ensemble is the least stable here, likely because its
Logistic-Regression meta-learner is itself being fit on a relatively small
number of out-of-fold base-learner predictions.

\section{Evaluation Metrics: ROC Curve and Confusion Matrices}
\begin{figure}[H]
\centering
\includegraphics[width=\linewidth]{figures/roc_curves.png}
\caption{ROC curves for tuned Bagging, Boosting (AdaBoost), and Stacked Ensemble.}
\label{fig:roc}
\end{figure}
\textbf{Inference:} Bagging achieves the highest test-set ROC-AUC (0.993),
narrowly ahead of the Stacked Ensemble (0.986) and Boosting (0.982). This
is a useful complement to the cross-validated accuracy ranking in
Section~9: Boosting is the most accurate at the default 0.5 threshold on
average, while Bagging ranks positive and negative cases slightly better
across all possible thresholds on this particular test split.

\begin{figure}[H]
\centering
\includegraphics[width=0.6\linewidth]{figures/cm_bagging.png}
\caption{Confusion matrix: tuned Bagging.}
\label{fig:cm_bag}
\end{figure}
\begin{figure}[H]
\centering
\includegraphics[width=0.6\linewidth]{figures/cm_boosting.png}
\caption{Confusion matrix: tuned Boosting (AdaBoost).}
\label{fig:cm_boost}
\end{figure}
\begin{figure}[H]
\centering
\includegraphics[width=0.6\linewidth]{figures/cm_stacking.png}
\caption{Confusion matrix: tuned Stacked Ensemble.}
\label{fig:cm_stack}
\end{figure}
\textbf{Inference:} All three tuned models make the same small number of
total errors (5 of 114 test samples) but distribute them slightly
differently between false positives and false negatives, which is why
their headline accuracy is identical while their ROC-AUC and CV stability
differ.

\section{Bias--Variance Analysis}

\begin{tabular}{lccc}
\toprule
Model & Train Accuracy & Validation Accuracy & Train--Val Gap \\
\midrule
Single Decision Tree      & 1.0000 & 0.9165 & 0.0835 \\
Bagging (Tuned)           & 0.9918 & 0.9604 & 0.0313 \\
Boosting (Tuned)          & 1.0000 & 0.9802 & 0.0198 \\
Stacked Ensemble (Tuned)  & 0.9648 & 0.9538 & \textbf{0.0110} \\
\bottomrule
\end{tabular}

\begin{figure}[H]
\centering
\includegraphics[width=0.8\linewidth]{figures/bias_variance_analysis.png}
\caption{Training vs. validation accuracy: single Decision Tree vs. all three ensemble methods.}
\label{fig:bias_variance}
\end{figure}
\textbf{Inference:} The train--validation gap shrinks steadily from the
single Decision Tree (8.35 points) through Bagging (3.13 points) and
Boosting (1.98 points) down to the Stacked Ensemble (1.10 points), directly
visualizing the variance-reduction promised by ensembling. Notably,
Boosting reaches a perfect 1.000 training accuracy (same as the single
tree) yet still generalizes far better (0.980 vs.\ 0.916 validation
accuracy) -- illustrating that boosting's sequential error-correction can
fit the training data extremely well \emph{without} the same overfitting
penalty a single unconstrained tree pays, because each weak learner in the
ensemble is individually simple (shallow), and the ensemble's effective
capacity is controlled by the learning rate and number of rounds rather
than by any single learner's depth.

\section{Observations}
\begin{itemize}
\item \textbf{How does Bagging reduce variance?} By training each base
tree on an independent bootstrap sample (and here, a random 50\% feature
subset as well), individual trees' errors become partially decorrelated.
Averaging over many decorrelated predictors reduces the ensemble's overall
variance roughly in proportion to how uncorrelated the individual errors
are, which is directly visible in Bagging's much smaller train--validation
gap (3.13 points) compared to the single Decision Tree (8.35 points), and
its lowest cross-validation standard deviation among the three ensembles
(0.0179).
\item \textbf{How does Boosting address model bias?} Each successive weak
learner in AdaBoost is explicitly trained to focus on the samples the
current ensemble misclassifies (via increased sample weights), so the
ensemble incrementally reduces the systematic errors (bias) that any single
simple learner would make. This lets Boosting achieve the highest
cross-validated accuracy of all three methods (98.02\%) using an ensemble
of otherwise-weak learners.
\item \textbf{Why does stacking benefit from heterogeneous models?} SVM
(margin-based), Naive Bayes (probabilistic, feature-independence
assumption), and Decision Tree (rule-based, axis-aligned splits) make
qualitatively different kinds of errors, since each imposes a different
inductive bias on the decision boundary. A meta-learner can, in principle,
learn to trust each base learner more in the regions of feature space where
it tends to be correct, capturing complementary strengths that no single
model family provides alone -- this is why dropping SVM from the base set
(the ``NB+DT only'' row of Table 3) reduced CV accuracy from 95.38\% to
94.29\%.
\item \textbf{Which ensemble method performed best and why?} By the most
reliable metric measured here -- 5-fold cross-validated accuracy -- tuned
\textbf{Boosting (AdaBoost, n\_estimators=100, learning\_rate=1.0)}
performed best (98.02\%), ahead of Bagging (96.04\%) and the Stacked
Ensemble (95.38\%). This is consistent with Boosting's bias-reduction
mechanism being particularly effective on this dataset: the underlying
classes are largely well-separated by a handful of features (as seen in
Experiment 5's feature-importance analysis), so a sequence of shallow,
sequentially-corrected weak learners can rapidly hone in on the true
decision boundary. Bagging remains competitive and is the most
\emph{stable} choice (lowest CV variance), which may be preferable in
settings where consistency across resamples of the data matters more than
squeezing out the last percentage point of accuracy.
\end{itemize}

\section{Conclusion}
Bagging, Boosting (AdaBoost and Gradient Boosting), and a Stacked Ensemble
(SVM + Naive Bayes + Decision Tree, with a Logistic Regression
meta-learner) were implemented and evaluated on the Wisconsin Diagnostic
Breast Cancer dataset using 5-fold cross-validation for hyperparameter
selection. All three tuned ensembles reached an identical 95.61\% single-
split test accuracy, but cross-validation revealed AdaBoost as the clear
overall winner (98.02\% CV accuracy) and Bagging as the most stable
(lowest fold-to-fold variance). The bias--variance analysis directly
confirmed the theoretical roles of each technique: Bagging cut the
train--validation gap from a single tree's 8.35 points down to 3.13 points
through variance reduction, while Boosting achieved an even smaller
generalization gap (1.98 points) despite reaching the same perfect
training accuracy as an unconstrained single tree, demonstrating bias
reduction through sequential error correction rather than simple averaging.
The experiment confirms that different ensemble strategies address
different failure modes of a single learner, and that the best choice
among them depends on which failure mode -- variance, bias, or model
diversity -- is most limiting for the dataset and task at hand.

\section{References}
\begin{enumerate}
\item Scikit-learn Developers, ``Ensemble methods,'' scikit-learn
documentation. [Online]. Available:
\url{https://scikit-learn.org/stable/modules/ensemble.html}
\item Scikit-learn Developers, ``BaggingClassifier,'' scikit-learn API
reference. [Online]. Available:
\url{https://scikit-learn.org/stable/modules/generated/sklearn.ensemble.BaggingClassifier.html}
\item W. H. Wolberg, W. N. Street, and O. L. Mangasarian, ``Breast Cancer
Wisconsin (Diagnostic) Data Set,'' UCI Machine Learning Repository, 1995.
[Online]. Available: \url{https://archive.ics.uci.edu}
\end{enumerate}

\end{document}
